\Needspace{9\baselineskip}\begin{center}\textit{Solving the indeterminate equation \(mxx+nyy=A\) by exclusions.}\end{center}


\Needspace{6\baselineskip}
\Article{323}
In Section V we showed how to find all representations of a given number \(A\) by the binary form \(mxx+nyy\),
\SourceLeaf{397}
or the solutions of the indeterminate equation \(mxx+nyy=A\), by a general method whose brevity, too, seems to leave nothing to be desired if all the values of the expression \(\surd{-mn}\) are already known modulo \(A\) itself and modulo that number divided by its square factors. Here, however, for the case where \(mn\) is positive, we shall explain a solution much more convenient than the direct one if those values must first be calculated for the latter. We shall assume that the numbers \(m\), \(n\) and \(A\) are positive and pairwise coprime, since the other cases can easily be reduced to this one. It is also evident that it is enough to find positive values of \(x,\allowbreak y\), since the others follow simply by changing the signs.

It is clear that \(x\) must be such that \(\frac{A-mxx}{n}\), for which we shall write \(V\), is positive, integral and a square. The first condition requires that \(x\) be no greater than \(\surd{\frac Am}\); the second holds of itself when \(n=1\), and otherwise requires that the value of the expression \(\frac Am\; (\mathrm{mod}.\ n)\) be a quadratic residue of \(n\). Denoting all the distinct values of the expression \(\surd{\frac Am}\; (\mathrm{mod}.\ n)\) by \(\pm r,\allowbreak \pm r'\), etc., we must have \(x\) of one of the forms \(nt+r,\allowbreak nt-r,\allowbreak nt+r'\), etc. The simplest procedure would therefore be to take all numbers of these forms below the bound \(\surd{\frac Am}\) (denoting their collection by \(\Omega\)), substitute them for \(x\), and retain only those for which \(V\) becomes a square. In the next article we shall show how to shorten this trial as much as we please.

\Needspace{6\baselineskip}
\Article{324}
The method of exclusion by which we shall do this consists, as in the preceding investigation, in taking several numbers at will, which here too we shall call \textit{exclusion numbers}, finding the values of \(x\) for which the value of \(V\) is a quadratic nonresidue of these exclusion numbers, and removing such values of \(x\) from \(\Omega\). Reasoning entirely analogous to that given in Article \(321\) shows that we need use only exclusion numbers that are primes or powers of primes. For an exclusion number of the latter kind, we need remove from the values of \(V\) only those of its nonresidues that are residues of all lower powers of the same prime, provided the exclusions with these have already been made.

Let the exclusion number therefore be \(E=p^\mu\), including the case \(\mu=1\), where \(p\) is a prime that does not divide \(m\), and suppose\HistoricalNote{*)}{For brevity, we include together the two cases in which \(n\) is and is not divisible by \(p\); in the latter we must put \(\nu=0\).} that \(p^\nu\) is the highest power
\SourceLeaf{398}
of the same prime by which \(n\) is divisible. Let \(a,\allowbreak b,\allowbreak c\), etc., be quadratic nonresidues of \(E\): all of them when \(\mu=1\), and the necessary ones, namely those that are residues of the lower powers, when \(\mu>1\). Calculate the roots of the congruences \(mz\equiv A-na,\allowbreak mz\equiv A-nb,\allowbreak mz\equiv A-nc\), etc., \((\mathrm{mod.}\ Ep^\nu=p^{\mu+\nu})\), and let them be \(\alpha,\allowbreak \beta,\allowbreak \gamma\), etc. It is easy to see that, if a value of \(x\) gives \(xx\equiv\alpha\) \((\mathrm{mod.}\ Ep^\nu)\), the corresponding value of \(V\) will be \(\equiv a\; (\mathrm{mod}.\ E)\), and thus a nonresidue of \(E\); and similarly for the other numbers \(\beta,\allowbreak \gamma\), etc. Conversely, it is equally easy to see that, if a value of \(x\) produces \(V\equiv a\; (\mathrm{mod}.\ E)\), then for the same value we have \(xx\equiv\alpha\; (\mathrm{mod}.\ Ep^\nu)\). Thus all values of \(x\) for which \(xx\) is congruent to none of the numbers \(\alpha,\allowbreak \beta,\allowbreak \gamma\), etc., modulo \(Ep^\nu\) produce values of \(V\) congruent to none of the numbers \(a,\allowbreak b,\allowbreak c\), etc., modulo \(E\). Now select from the numbers \(\alpha,\allowbreak \beta,\allowbreak \gamma\), etc., all the quadratic residues of \(Ep^\nu\), and call them \(g,\allowbreak g',\allowbreak g''\), etc. Calculate the values of the expressions \(\surd g,\allowbreak \surd{g'}.\surd{g''}\), etc., \((\mathrm{mod.}\ Ep^\nu)\), and suppose that these give \(\pm h,\allowbreak \pm h',\allowbreak \pm h''\), etc. Once this has been done, it is clear that all numbers of the forms \(Ep^\nu t\pm h,\allowbreak Ep^\nu t\pm h',\allowbreak Ep^\nu t\pm h''\), etc., can safely be removed from \(\Omega\), and that no value of \(x\) remaining in \(\Omega\) after this exclusion can correspond to a value of \(V\) of the forms \(Eu+a,\allowbreak Eu+b,\allowbreak Eu+c\), etc. Moreover, it is evident that such values of \(V\) cannot arise from any value of \(x\) at all if none of the numbers \(\alpha,\allowbreak \beta,\allowbreak \gamma\), etc., are quadratic residues of \(Ep^\nu\); hence in this case the number \(E\) cannot be used as an exclusion number. — We may use as many exclusion numbers of this kind as we please, and thus reduce the collection of numbers in \(\Omega\) as much as we please.

Let us now see whether primes dividing \(m\), or powers of such primes, can also be used as exclusion numbers. Let \(B\) be the value of the expression \(\frac An\; (\mathrm{mod}.\ m)\). It is clear that \(V\) is always congruent to \(B\) modulo \(m\), whatever value is taken for \(x\), and thus that a necessary condition for the proposed equation to be solvable is that \(B\) be a quadratic residue of \(m\). Now let \(p\) denote any odd prime divisor of \(m\), which by hypothesis divides neither \(n\) nor \(A\), and hence does not divide \(B\) either. For any value of \(x\), \(V\) will be a residue of \(p\), and therefore of any power of \(p\) as well; hence \(p\) and its powers cannot serve as exclusion numbers. By exactly similar reasoning, when \(m\) is divisible by \(8\), a necessary condition for the proposed equation to be solvable is \(B\equiv1\; (\mathrm{mod}.\ 8)\). Then \(V\), for any value of \(x\), will also be \(\equiv1\) \((\mathrm{mod.}\ 8)\), and powers of two are therefore unsuitable for exclusion. — But
\SourceLeaf{399}
when \(m\) is divisible by \(4\) but not by \(8\), similar reasoning requires \(B\equiv1\) \((\mathrm{mod.}\ 4)\), and hence the value of the expression \(\frac An\; (\mathrm{mod}.\ 8)\) is either \(1\) or \(5\); denote it by \(C\). It is easy to see that, for an even value of \(x\), we have \(V\equiv C\), and for an odd value, \(V\equiv C+4\; (\mathrm{mod}.\ 8)\). Thus even values must be rejected when \(C=5\), and odd values when \(C=1\). — Finally, when \(m\) is divisible by \(2\) but not by \(4\), let \(C\), as before, be the value of the expression \(\frac An\; (\mathrm{mod}.\ 8)\), which will be \(1,\allowbreak 3,\allowbreak 5\) or \(7\); and let \(D\) be the value of \(\frac{\frac12m}{n}\; (\mathrm{mod}.\ 4)\), which will be \(1\) or \(3\). Now the value of \(V\) is evidently always \(\equiv C-2Dxx\; (\mathrm{mod}.\ 8)\), and hence, for even \(x\), \(\equiv C\), and for odd, \(\equiv C-2D\). It follows easily that all odd values of \(x\) must be rejected when \(C=1\); all even values when \(C=3\) and \(D=1\), or \(C=7\) and \(D=3\); and that all remaining values produce \(V\equiv1\; (\mathrm{mod}.\ 8)\), a residue of any power of two. In the other cases, namely when \(C=5\), or \(C=3\) and \(D=3\), or \(C=7\) and \(D=1\), we have \(V\equiv3,\allowbreak 5\) or \(7\; (\mathrm{mod}.\ 8)\), whether \(x\) is taken even or odd. Thus in these cases the proposed equation has no solution at all.

Moreover, just as we have shown here how to find the value of \(x\) by exclusions, we can, with the necessary changes, obtain the value of \(y\). The method of exclusion can therefore always be applied to the proposed problem in two ways, unless \(m=n=1\), when they coincide. Of these, the one for which \(\Omega\) contains fewer terms is generally preferable, and this can easily be estimated beforehand. — Finally, we hardly need to observe that, if after some exclusions \textit{all} numbers have disappeared from \(\Omega\), this must be regarded as a certain indication that the proposed equation is impossible.

\Needspace{6\baselineskip}
\Article{325}
\textit{Example.} Let the equation \(3xx+455yy=10857362\) be given. We shall solve it in two ways, \textit{first} by finding the values of \(x\), and then those of \(y\). The bound for the former in this case is \(\surd{3619120\frac23}\), which lies between \(1902\) and \(1903\); the value of the expression \(\frac A3\; (\mathrm{mod}.\ 455)\) is \(354\), and the values of the expression \(\surd{354}\; (\mathrm{mod}.\ 455)\) are \(\pm82,\allowbreak \pm152\), \(\pm173,\allowbreak \pm212\). Hence \(\Omega\) consists of the following \(33\) numbers: \(82,\allowbreak 152,\allowbreak 173,\allowbreak 212,\allowbreak \) \(243,\allowbreak 282,\allowbreak 303,\allowbreak 373,\allowbreak 537,\allowbreak 607,\allowbreak 628,\allowbreak 667,\allowbreak 698,\allowbreak 737,\allowbreak 758,\allowbreak 828,\allowbreak 992,\allowbreak 1062,\allowbreak 1083,\allowbreak \) \(1122,\allowbreak 1153,\allowbreak 1192,\allowbreak 1213,\allowbreak 1283,\allowbreak 1447,\allowbreak 1517,\allowbreak 1538,\allowbreak 1577,\allowbreak 1608,\allowbreak 1647,\allowbreak 1668,\allowbreak 1738,\allowbreak \) \(1902\). The number \(3\) cannot be used for exclusion in this case, because it divides \(m\). For the exclusion number \(4\) we have \(a=2,\allowbreak b=3\), giving \(\alpha=0,\allowbreak \beta=3\);
\SourceLeaf{400}
\(g=0\), and the values of the expression \(\surd{g}\ (\mathrm{mod.}\ 4)\) are \(0\) and \(2\). Thus all numbers of the forms \(4t\) and \(4t+2\), \textit{i.e.} all even numbers, must be removed from \(\Omega\); denote the sixteen remaining numbers by \(\Omega'\). For \(E=5\), which also divides \(n\), the congruences \(mz\equiv A-2n\) and \(mz\equiv A-3n\ (\mathrm{mod.}\ 25)\) have roots \(9\) and \(24\), both residues of \(25\), and the values of the expressions \(\surd9\) and \(\surd{24}\ (\mathrm{mod.}\ 25)\) are \(\pm3,\allowbreak \ \pm7\). Removing from \(\Omega'\) all numbers of the forms \(25t\pm3,\allowbreak \ 25t\pm7\) leaves these ten \((\Omega'')\): \(173,\allowbreak 373,\allowbreak 537,\allowbreak 667,\allowbreak 737,\allowbreak 1083,\allowbreak 1213,\allowbreak 1283,\allowbreak 1517,\allowbreak 1577\). For \(E=7\), the congruences \(mz\equiv A-3n,\allowbreak \ mz\equiv A-5n,\allowbreak \ mz\equiv A-6n\) \((\mathrm{mod.}\ 49)\) have roots \(32,\allowbreak 39,\allowbreak 18\), all residues of \(49\), and the values of the expressions \(\surd{32},\allowbreak \surd{39},\allowbreak \surd{18}\ (\mathrm{mod.}\ 49)\) are \(\pm9,\allowbreak \pm23,\allowbreak \pm19\). Removing from \(\Omega''\) the numbers of the forms \(49t\pm9,\allowbreak \ 49t\pm19,\allowbreak \ 49t\pm23\) leaves these five \((\Omega''')\): \(537,\allowbreak 737\), \(1083,\allowbreak 1213,\allowbreak 1517\). For \(E=8\) we have \(a=5\), giving \(\alpha=5\), a nonresidue of \(8\); hence the exclusion number \(8\) cannot be used. The number \(9\) must be skipped for the same reason as \(3\). For \(E=11\), the numbers \(a,\allowbreak b\), etc., are \(2,\allowbreak 6\), \(7,\allowbreak 8,\allowbreak 10\); \(\nu=0\); hence \(\alpha,\allowbreak \beta\), etc., \(=8,\allowbreak 10,\allowbreak 5,\allowbreak 0,\allowbreak 1\), of which three are residues of \(11\), namely \(0,\allowbreak 1,\allowbreak 5\). It follows that we must remove from \(\Omega'''\) the numbers of the forms \(11t,\allowbreak \ 11t\pm1,\allowbreak \ 11t\pm4\), leaving \(537,\allowbreak 1083,\allowbreak 1213\). Trying these gives for \(V\) the respective values \(21961,\allowbreak 16129,\allowbreak 14161\), of which only the second and third are squares. Thus the proposed equation admits two solutions with positive values of \(x,\allowbreak y\), namely \(x=1083,\allowbreak y=127\) and \(x=1213,\allowbreak y=119\).

\textit{Second.} If we wish to find the other unknown of the same equation by exclusions, put the equation in the form \(455xx+3yy=10857362\), interchanging \(x\) and \(y\) so that all the notation of Articles \(323,\allowbreak 324\) may be retained. Here the bound for the values of \(x\) lies between \(154\) and \(155\); the value of the expression \(\frac Am\ (\mathrm{mod.}\ n)\) is \(1\), and the values of \(\surd1\ (\mathrm{mod.}\ 3)\) are \(+1\) and \(-1\). Hence \(\Omega\) contains all numbers of the forms \(3t+1\) and \(3t-1\), \textit{i.e.} all numbers not divisible by \(3\) up to and including \(154\), whose number is \(103\). Applying the rules given above, however, we find that for the exclusion numbers \(3;\allowbreak \ 4;\allowbreak \ 9;\allowbreak \ 11;\allowbreak \) \(17;\allowbreak \ 19;\allowbreak \ 23\) we must reject the numbers of the forms \(9t\pm4;\allowbreak \ 4t,\allowbreak \ 4t\pm2\), or all even numbers; \(27t\pm1,\allowbreak \ 27t\pm10;\allowbreak \ 11t,\allowbreak \ 11t\pm1,\allowbreak \ 11t\pm3;\allowbreak \ 17t\pm3,\allowbreak \ 17t\pm4,\allowbreak \ 17t\pm5\), \(17t\pm7;\allowbreak \ 19t\pm2,\allowbreak \ 19t\pm3,\allowbreak \ 19t\pm8,\allowbreak \ 19t\pm9;\allowbreak \ 23t,\allowbreak \ 23t\pm1,\allowbreak \ 23t\pm5,\allowbreak \ 23t\pm7\), \(23t\pm9,\allowbreak \ 23t\pm10\). Removing these leaves \(119,\allowbreak 127\), the only two that give \(V\) a square value, yielding the same solutions as those obtained above.
\SourceLeaf{401}
